\section{Parting Ways: Jewish Withdrawal and the New Greek Versions}

As the Church made the Greek Bible its own, Greek-speaking Jewry moved the
other way---not by a single decree (none is attested), but along the path
the Na\d{h}al \d{H}ever scroll had already marked: toward the Hebrew.

The polemical friction is audible by the mid-second century. Justin tells
Trypho that the Jewish teachers ``refuse to admit that the interpretation
made by the seventy elders \ldots\ is a correct one'' and ``attempt to
frame another,'' and he accuses them of deleting messianic passages---among
them the words ``from the wood'' in Psalm 96's ``the Lord reigned.''
Modern scholarship has largely reversed the charge: the disputed words are
absent from the Hebrew and from the best Greek witnesses alike, and most of
Justin's ``deletions'' are Christian expansions that his copies carried,
not Jewish excisions. The accusation is nonetheless first-rate evidence
for the situation: by c.~160 the two communities were reading measurably
different Greek Bibles and each blamed the other.\endnote{Justin,
\work{Dialogue with Trypho} 71--73 (ANF vol.~1), checked at New Advent;
the modern assessment of the ``removed'' passages follows the standard
critical literature (class~K).}

Three named versions replaced the old translation in Jewish use, all known
today almost entirely from fragments and from Origen's and the Fathers'
reports.

\textbf{Aquila.} A rigorously literal version---calqued Hebrew syntax,
etymologizing renderings---attested by Irenaeus already as the work of
Aquila of Pontus, a Jewish proselyte. Epiphanius, writing two centuries
later, dates him under Hadrian (early second century) and makes him a
relative of the emperor, a lapsed Christian, and a convert to Judaism; the
biographical detail is late and partly hostile, but the dating fits
Irenaeus's knowledge of him. Rabbinic tradition, strikingly, remembers the
same figure with honor: in the Jerusalem Talmud, Aquila the proselyte
translates the Torah before Rabbis Eliezer and Joshua, who praise him in
the language of Psalm 45's tribute to beauty.\endnote{Irenaeus, \work{Against Heresies} 3.21.1; Epiphanius,
\work{On Weights and Measures} 14--15 (Dean translation, paraphrased);
Jerusalem Talmud, \work{Megillah} 1:9 (Guggenheimer translation at
Sefaria, paraphrased). The rabbinic praise shows Aquila's version was a
Jewish project with rabbinic approbation, not merely an anti-Christian
maneuver.}

\textbf{Symmachus.} A version in idiomatic, literary Greek, faithful to
the Hebrew sense without Aquila's violence to Greek. Eusebius calls its
author an Ebionite (a Jewish Christian); Epiphanius makes him a Samaritan
convert to Judaism; the identifications cannot both be right and neither
can be verified.\endnote{Eusebius, \work{Church History} 6.17 (NPNF
series~2), checked at New Advent; Epiphanius as reported in the modern
surveys. The disagreement is preserved, not resolved.}

\textbf{Theodotion---and the Theodotion problem.} A revision of the old
Greek toward the Hebrew, traditionally placed in the later second century.
The difficulty, recognized since antiquity and sharpened by Barth\'elemy,
is that characteristically ``Theodotionic'' readings appear in texts
older than Theodotion---including the New Testament---and that the
church's preferred Greek Daniel, transmitted under his name, displaced the
Old Greek Daniel almost everywhere. Either ``Theodotion'' finished a
revising tradition (\textit{kaige}) generations older than himself, or
the name shelters more than one enterprise; the question remains
open.\endnote{Irenaeus already names Theodotion (\work{Against Heresies}
3.21.1) a generation after Justin quotes ``Theodotionic'' wording; the
proto-Theodotion/\textit{kaige} analysis is Barth\'elemy's (class~K),
refined and contested since. The near-total loss of Old Greek Daniel
from the manuscript tradition is uncontested.}

The withdrawal was gradual, regional, and never total---Greek-speaking
Jews were still hearing Greek Scripture centuries later---but its
direction was settled: the version the Church had adopted ceased to be a
Jewish book, and the mourning notice in Soferim shows how some rabbinic
memory finally judged the whole episode. As late as 553, Justinian's
Novella 146 was still regulating a live dispute: it permitted synagogue
reading in Greek---commending the Septuagint, allowing Aquila---against
those who insisted on Hebrew alone.\endnote{\work{Soferim} 1.7, cited
above; Justinian, Novella 146 \work{De Hebraeis} (553), consulted in the
Blume translation published by the University of Wyoming law library. The
novella attests continuing Jewish Greek-Bible use in the sixth century,
imperial interference in it, and a rabbinic preference for Hebrew---each
within the law's own polemical frame.}
